\section*{الفصل \ref{ch:b3:exploration-seismology} --- الزلزالية الاستكشافية}

\begin{solution}{exo:b3:exploration-seismology:1}
$Z = \rho v$: الطباشير $2000\times2500 = \qty{5.00e6}{kg.m^{-2}.s^{-1}}$، الطَّفل
$2300\times2300 = \qty{5.29e6}{kg.m^{-2}.s^{-1}}$، الحجر الجيري $2600\times4500 =
\qty{11.70e6}{kg.m^{-2}.s^{-1}}$. الطباشير--الطَّفل: $R = (5.29 - 5.00)/(5.29 + 5.00)
= 0.29/10.29 = 0.0282$؛ الطَّفل--الحجر الجيري: $R = (11.70 - 5.29)/(11.70 + 5.29) =
6.41/16.99 = 0.377$. الطاقة المنعكسة هي $R^2 = 0.142$: أعلى
الحجر الجيري يُعيد \qty{14.2}{\percent} من الطاقة الساقطة، أما السطح الفاصل بين الطباشير--الطَّفل
فيُعيد أقل من \qty{0.1}{\percent}.
\end{solution}

\begin{solution}{exo:b3:exploration-seismology:2}
$t_0 = 2\times800/2000 = \qty{0.800}{s}$؛ $t = \sqrt{0.800^2 + (1000/2000)^2} =
\sqrt{0.89} = \qty{0.9434}{s}$؛ الانزياح الزمني هو \qty{0.1434}{s}. التقريب
يعطي $1000^2/(2\times2000^2\times0.800) = \qty{0.1563}{s}$، أكبر
بمقدار \qty{9}{\percent}: عند إزاحة أكبر من العمق، لم
يعد الشرط $x \ll vt_0$ صحيحًا.
\end{solution}

\begin{solution}{exo:b3:exploration-seismology:3}
تنتقل نقاط المنتصف بمقدار نصف فترة المستقبِلات، فتكون فترة \omterm{def:b3:exploration-seismology:cmp}{نقطة المنتصف المشتركة}
هي $12.5/2 = \qty{6.25}{m}$. \omterm{def:b3:exploration-seismology:cmp}{تعدّدية التغطية} هي $240\times12.5/(2\times25) = 60$. أما تعدّدية
120 تحتاج إلى طلقة كل $240\times12.5/(2\times120) = \qty{12.5}{m}$.
\end{solution}

\begin{solution}{exo:b3:exploration-seismology:4}
$\lambda = v/f$: \qty{100}{m} في الحجر الرملي، $\lambda/4 = \qty{25}{m}$؛
\qty{150}{m} في الحجر الجيري، $\lambda/4 = \qty{37.5}{m}$. مع
$\lambda = 2800/30 = \qty{93.3}{m}$ و$z = \qty{1000}{m}$،
$r_F \approx \sqrt{\lambda z/2} = \sqrt{46\,667} = \qty{216}{m}$ (وتصبح \qty{217}{m}
مع الحد $\lambda^2/16$).
\end{solution}

\begin{solution}{exo:b3:exploration-seismology:5}
\omterm{thm:b3:exploration-seismology:dix}{معادلة ديكس}: $v^2 = (2800^2\times0.700 - 2400^2\times0.500)/0.200 = (5\,488\,000 -
2\,880\,000)/0.200 = \qty{1.304e7}{m^2.s^{-2}}$، إذن $v = \qty{3611}{m.s^{-1}}$.
تُجتاز الطبقة في \qty{0.100}{s} باتجاه واحد: $h = 3611\times0.100 =
\qty{361}{m}$.
\end{solution}

\begin{solution}{exo:b3:exploration-seismology:6}
(أ) ليكن $d$ المسافة من نقطة المنتصف M إلى العاكس، بحيث
$t_0 = 2d/v$. الشعاع المنعكس من المصدر S إلى المستقبِل G طولُه
يساوي طول الخط المستقيم من الصورة S$'$ للمصدر S في العاكس إلى G. نضع M عند
مبدأ الإحداثيات، والخط على امتداد الانحدار؛ وتقع S$'$ على بُعد $2d_S$ من S على طول
العمودي، مع $d_S = d - (x/2)\sin\varphi$. مربّع طول S$'$G هو
$x^2 + 4d_S^2 + 4d_Sx\sin\varphi = x^2 + 4d^2 - x^2\sin^2\varphi$، جميع الحدود
التي تتضمّن $dx$ فتتلاشى، إذن $v^2t^2 = 4d^2 + x^2\cos^2\varphi$
و$t^2 = t_0^2 + x^2\cos^2\varphi/v^2$. (ب) تتبع المجموعة قطعًا زائديًا سرعته
$V = v/\cos\varphi = 2000/\cos\ang{12} = \qty{2045}{m.s^{-1}}$: الانحدار يجعل
\omterm{def:b3:exploration-seismology:velocities}{سرعة التكديس} أعلى من اللازم، هنا بمقدار \qty{2.2}{\percent}.
\end{solution}

\begin{solution}{exo:b3:exploration-seismology:7}
بحساب $\sqrt{60} = 7.75$. تتطلّب مضاعفة نسبة الإشارة إلى الضجيج أربعة أضعاف
عدد المسارات، أي تعدّدية 240.
\end{solution}

\begin{solution}{exo:b3:exploration-seismology:8}
يتغيّر العمق المرسوم $vt/2$ بمقدار $\qty{1.5}{km.s^{-1}}\times\qty{0.40}{s}
= \qty{0.60}{km}$ لكل كيلومتر: $\tan\theta_a = 0.60$، $\theta_a = \ang{31.0}$.
بمبرهنة الإزاحة الصفرية $\sin\theta = 0.60$: $\theta = \ang{36.9}$.
\end{solution}

\begin{solution}{exo:b3:exploration-seismology:9}
يضع كل نتوء نسخة من $w$ متمركزة عليه: $0.3\,w$ على العيّنات 0--2 يعطي
$(-0.15, 0.3, -0.15)$، و$-0.2\,w$ على العيّنات 2--4 يعطي $(0.1, -0.2, 0.1)$.
ومجموعهما هو $(0, -0.15, 0.3, -0.05, -0.2, 0.1, 0)$، مكتوبًا من عيّنة واحدة
قبل المتسلسلة إلى عيّنة واحدة بعدها. تبلغ الانعكاسات ذروتها عند العيّنتين 1 و3، بإشاراتها
الحقيقية؛ وتتراكب عند العيّنة 2، حيث $-0.15 + 0.1 = -0.05$: تداخلٌ
ضعيف، لكنّ الذروتين متمايزتان.
\end{solution}

\begin{solution}{exo:b3:exploration-seismology:10}
$\arcsin(1800/2900) = \ang{38.37}$ و$\arcsin(1800/4400) = \ang{24.15}$، إذن
$i_c = (38.37 + 24.15)/2 = \ang{31.26}$ و$\theta = (38.37 - 24.15)/2 =
\ang{7.11}$؛ $v_2 = 1800/\sin\ang{31.26} = \qty{3469}{m.s^{-1}}$. المتوسط
التوافقي $2\times2900\times4400/(2900 + 4400) = \qty{3496}{m.s^{-1}}$ أعلى
بمقدار \qty{0.8}{\percent}؛ وبضربه في $\cos\theta = 0.9923$ يعطي
\qty{3469}{m.s^{-1}}، وهو التقريب المعتاد للانحدار الصغير
$v_2 \approx \cos\theta\,2v_dv_u/(v_d + v_u)$.
\end{solution}

\begin{solution}{exo:b3:exploration-seismology:11}
(أ) الموجة المركّبة هي $R[w(t) - w(t - \Delta t)] \approx R\,\Delta t\,w'(t)$ بواسطة
نشر تايلور من الرتبة الأولى. (ب) سعتها متناسبة مع
$\Delta t$، لذا فإن طبقة أرقّ من \omterm{def:b3:exploration-seismology:resolution}{سمك التوليف} ما تزال تُعيد إشارة
بشكل $w'$، لكن قمّتها وقاعدتها لا يمكن تمييزهما كحدثين منفصلين: إنّها
تُكتشف ولا تُفصَل. وإذا عُرفت $R$ والموجة، مثلًا من
بئرٍ، فإنّ السعة تعطي $\Delta t$ وبالتالي السُّمك. (ج) السعة متناسبة
مع السُّمك: النسبة هي $5/10 = 1/2$.
\end{solution}

\begin{solution}{exo:b3:exploration-seismology:12}
(أ) $t(x) = (2/v)\sqrt{z^2 + x^2}$: $t(0) = 2\times600/3000 = \qty{0.400}{s}$
و$t(300) = (2/3000)\sqrt{600^2 + 300^2} = \qty{0.4472}{s}$. (ب) $t^2 = t_0^2 +
x^2/(v/2)^2$، وهي صيغة \omterm{thm:b3:exploration-seismology:nmo}{الانزياح الزمني العادي} بالسرعة $v/2$. لو كان لانعكاسٍ زمنُ القمّة نفسه
في مقطع الإزاحة الصفرية، لما ظهر أيّ انزياح زمني؛ وبالمقارنة مع
قطعٍ زائد للانزياح العادي ذي $t_0$ نفسه في المجموعة، فإنّ الانزياح الزمني للحيود
$x^2/(2(v/2)^2t_0)$ أكبر أربع مرّات، لأنّ الزمن يعتمد على
المسافة من المبعثِر مرّتين، ذهابًا وإيابًا. (ج) يجمع الترحيل
المقطع على طول هذا القطع الزائد ويضع الطاقة عند قمّته: ينكمش
الحيود إلى نقطة عند $x = 0$، $z = \qty{600}{m}$.
\end{solution}

\begin{solution}{pb:b3:exploration-seismology:1}
\textbf{1.} من W، $400/(486.1 - 425.2) = 400/0.0609 = \qty{6568}{m.s^{-1}}$؛
من E، $400/(486.1 - 372.4) = 400/0.1137 = \qty{3518}{m.s^{-1}}$. موجة الرأس
القادمة من W تبدو سريعة: إنّ W هي النهاية السفلى للانحدار، وطلقاتها تنتقل صعودًا بعكس الانحدار؛
يرتفع السطح الانكساري من W إلى E.

\textbf{2.} $\arcsin(2400/3518) = \ang{43.02}$, $\arcsin(2400/6568) =
\ang{21.43}$: $i_c = \ang{32.22}$, $\theta = \ang{10.79}$.

\textbf{3.} $v_2 = 2400/\sin\ang{32.22} = \qty{4501}{m.s^{-1}}$: ضمن
مدى الحجر الجيري، أي قمّة الحجر الجيري الجوراسي.

\textbf{4.} أزمنة القطع هي $t - x/v_{\text{ظاهري}}$ عند \qty{1200}{m}:
$t_{i,W} = 0.4861 - 1200/6568 = \qty{0.3034}{s}$ و$t_{i,E} = 0.4861 -
1200/3518 = \qty{0.1450}{s}$. مع $\cos i_c = 0.8460$، $h = v_1t_i/(2\cos
i_c)$: $h_W = 2400\times0.3034/1.692 = \qty{430.4}{m}$ و$h_E =
\qty{205.7}{m}$.

\textbf{5.} بقسمة على $\cos\theta = 0.9823$: \qty{438.1}{m} تحت W
و\qty{209.4}{m} تحت E. ثمّ $(438.1 - 209.4)/1200 = 0.1906 =
\tan\ang{10.79}$، وهو انحدار السؤال 2.

\textbf{6.} لا يُنتج الطَّفل موجة رأس: فهو غير مرئي، ويُدرج المسح
الطَّفل مع الطباشير. هنا أُعطيت $v_1 = \qty{2.40}{km.s^{-1}}$ بوصفها
متوسّط الاثنين؛ فلو اعتُمدت سرعة الموجة المباشرة للطباشير، \qty{2.50}{km.s^{-1}}،
على الغطاء الفوقي كلّه، لخرجت كلّ الأعماق أكبر من الحقيقة.

\textbf{7.} $V = 600/\sqrt{0.5437^2 - 0.4878^2} = 600/0.2401 =
\qty{2499}{m.s^{-1}}$ و$V = 600/\sqrt{0.6240^2 - 0.5748^2} =
\qty{2471}{m.s^{-1}}$. القيمة الثانية تأخذ متوسّط الطباشير مع الطَّفل الأبطأ
الذي تحته، فينخفض المتوسّط.

\textbf{8.} \qty{2811}{m.s^{-1}} و\qty{2829}{m.s^{-1}}.

\textbf{9.} الطبقة 1: \qty{2499}{m.s^{-1}}. الطبقة 2: $v^2 = (2471^2\times0.5748
- 2499^2\times0.4878)/0.0870$، $v = \qty{2308}{m.s^{-1}}$. الطبقة 3:
$v^2 = (2811^2\times0.6548 - 2471^2\times0.5748)/0.0800$,
$v = \qty{4561}{m.s^{-1}}$. الطبقة 4: $v^2 = (2829^2\times0.7348 -
2811^2\times0.6548)/0.0800$, $v = \qty{2972}{m.s^{-1}}$.

\textbf{10.} الطباشير (\qty{2499}{m.s^{-1}})، الطَّفل (\qty{2308}{m.s^{-1}})،
الحجر الجيري (\qty{4561}{m.s^{-1}})، الحجر الرملي (\qty{2972}{m.s^{-1}})، كلّ منها ضمن
مداها المنشور. الانعكاسات هي سطح التماس بين الطباشير--الطَّفل، وقمّة
الحجر الجيري، وسطح التماس بين الحجر الجيري--الحجر الرملي، وعدم التوافق عند
قمّة قاعدة الشيست.

\textbf{11.} الطباشير والطَّفل لهما المعاوقة نفسها تقريبًا: $R = 0.028$
(التمرين 1)، مقابل 0.38 لقمّة الحجر الجيري.

\textbf{12.} القطع الزائد للسرعة $V_{\text{rms}}$ يصحّ فقط حتى الرتبة الثانية
في الإزاحة؛ وعند إزاحات مقدارها \qty{2}{km}، أي ضعف عمق أعمق
عاكس، تنحرف الأزمنة عنه، ولو طُوبق قطعٌ زائد على الانتشار كلّه
لأعطى سرعات خاطئة، والفترات المحسوبة بطريقة ديكس أسوأ من ذلك. وتُحدث الإزاحات البعيدة
مطًّا في الموجات القصيرة عند تصحيحها.

\textbf{13.} أزمنة الذهاب والإياب في الطبقات: \qty{0.4878}{s}، \qty{0.0870}{s}،
\qty{0.0800}{s}، \qty{0.0800}{s}. السماكات $v\,\Delta t/2$: \qty{609.5}{m}،
\qty{100.4}{m}, \qty{182.4}{m}, \qty{118.9}{m}.

\textbf{14.} الأعماق تحت السطح: \qty{609.5}{m}، \qty{709.9}{m}،
\qty{892.3}{m}، \qty{1011.2}{m}. جاءت قيم الطَّفل والحجر الجيري والحجر الرملي
ضمن \qty{0.4}{m}، \qty{2.4}{m} و\qty{1.1}{m} من سماكاتها، ويزداد
الخطأ في الحجر الجيري السريع الرقيق، حيث يُكلّف خطأ صغير في الزمن
أكثر ما يكون.

\textbf{15.} $\bar v = 2\times1011.2/0.7348 = \qty{2752}{m.s^{-1}}$. سرعة RMS
تُرجّح كلّ طبقة بمقدار $v^2$، ولا تقلّ أبدًا عن المتوسّط المرجَّح
بالزمن: $2829\times0.7348/2 = \qty{1039}{m}$، \qty{28}{m} أعمق مما ينبغي.

\textbf{16.} $\lambda = 2308/30 = \qty{76.9}{m}$, $\lambda/4 = \qty{19.2}{m}$;
وتقدير \omterm{def:b3:exploration-seismology:wavelet}{موجة ريكر القصيرة} هو $2308/(5.2\times30) = \qty{14.8}{m}$. الطَّفل
سماكته \qty{100}{m}، أي خمسة أمثال \omterm{def:b3:exploration-seismology:resolution}{سمك التوليف}: قمّته وقاعدته
مُحلَّلتان. ويتوالف الطَّفل، بحسب قاعدة ربع الطول الموجي، عند $f = v/(4b) =
2308/400 = \qty{5.77}{Hz}$.

\textbf{17.} $\lambda = 2829/30 = \qty{94.3}{m}$, $z = \qty{1011}{m}$:
$r_F \approx \sqrt{94.3\times1011/2} = \qty{218}{m}$ (أو \qty{220}{m} مع
الحدّ $\lambda^2/16$).

\textbf{18.} $t(x)^2 = t_0^2 + (2x/v)^2$، إذن $v = 2x/\sqrt{t^2 - t_0^2} =
400/\sqrt{0.5732^2 - 0.550^2} = 400/0.1614 = \qty{2478}{m.s^{-1}}$.

\textbf{19.} $z = vt_0/2 = 2478\times0.550/2 = \qty{681}{m}$.

\textbf{20.} يتغيّر العمق المرسوم $vt/2$ بمقدار $1.24\times0.167 = 0.207$ لكلّ
وحدة طول: $\theta_a = \arctan0.207 = \ang{11.70}$. وللانحدار الحقيقي
$\sin\theta = 0.207$: $\theta = \ang{11.95}$، وهو انحدار الجناح.

\textbf{21.} $d = 2480\times0.655/2 = \qty{812}{m}$؛ يُزيح الترحيل
الانعكاس $d\sin\theta = 812\times0.207 = \qty{168}{m}$ صعودًا بعكس الانحدار، نحو
قمّة الطية المحدبة، ويرتفع من \qty{812}{m} إلى $d\cos\theta =
\qty{795}{m}$.

\textbf{22.} $2500\times0.160/2 = \qty{200}{m}$.

\textbf{23.} من السؤالين 13 و14، يقع عدم التوافق القاعدي
\textbf{\qty[text-series-to-math=true]{1011}{m}} تحت السطح عند \babelsublr{CMP S}.
\end{solution}
